\section{Introduction}

Tensor networks provide a structured representation for high-dimensional arrays and are widely used in quantum many-body simulation, statistical physics, graphical-model inference, and tensor algorithms~\cite{orus2014practical}. A large body of work on tensor network contraction (TNC) has focused on three recurring questions: how to design better contraction orders, how to control the size of intermediate states, and how to reduce the cost of approximate contraction by compressing those intermediate states~\cite{markov2008simulating,dumitrescu2018benchmarking,jakes2019carving,gray2021hyperoptimized,huang2021efficient,pan2020contracting,gray2024hyperoptimized,ma2024approximate}. Most of this literature targets scalar-valued contraction, local observables, or local queries in probabilistic inference, rather than explicitly generating the full tensor that remains after preserving open legs~\cite{orus2014practical,dechter1998bucket,pan2020contracting,heddes2026approximating}. Once full materialization is required, both compute and memory costs can increase dramatically~\cite{jakes2019carving}.

This paper studies precisely that regime: tensor network contraction with open legs when the final output must be emitted as a dense tensor. The key distinction from standard approximate contraction is that the algorithm cannot change the form of the result. It cannot replace the final answer with a compressed representation, nor can it return only a few partial observables. Instead, the goal is to reduce the cost of producing the same dense output. Writing the end-to-end time as
\[
T_{\mathrm{total}} = T_{\mathrm{contract}} + T_{\mathrm{emit}},
\]
we interpret $T_{\mathrm{contract}}$ as the internal cost of contraction, merging, and intermediate-state processing before the final tensor is written out, and $T_{\mathrm{emit}}$ as the cost of materializing those results in the final output buffer. Because the dense output itself is mandatory, $T_{\mathrm{emit}}$ is usually unavoidable. The main algorithmic target is therefore $T_{\mathrm{contract}}$.

Motivated by this observation, we develop an approximate contraction framework for full materialization. The central idea is to build a partition tree over the interaction graph, recursively maintain separator states for sub-networks, and keep the output boundary explicit throughout the computation. Approximation is applied only to separator states that propagate inside the network. In other words, we do not approximate the final output object itself; we approximate the internal states that must be traversed in order to produce that output. This perspective is what distinguishes our setting from methods that directly compress the output side and from methods designed for scalar contraction.

From a representation viewpoint, each subtree state can be interpreted as an interface matrix from output-side configurations to boundary-side configurations. We propagate these states on the tree in low-rank factored form and recompress them after each merge to control rank growth. To avoid the common bottleneck in which interface dimensions inflate before being compressed again, we further introduce implicit sketch-based merging: rather than explicitly forming the merged interface, we recover its dominant subspace through operator access. The result is a unified pipeline that combines partition-tree recursion, separator-state approximation, and sketch-driven recompression while preserving output semantics.

Our experiments address three questions. First, can \method{} provide a practically useful accuracy-time trade-off for full-materialization TNC? Second, do the observed gains genuinely come from the contraction stage rather than from output write-back artifacts? Third, how do approximation strength and key mechanisms affect behavior? On representative classical networks, \method{} keeps the relative error below $4.81\times 10^{-3}$ while achieving speedups from $13.98\times$ to $84.05\times$ over exact TNC. A fixed small-rank baseline is faster in some cases, but its relative error rises to $0.65$--$0.98$. Additional analysis shows that the end-to-end gains are almost entirely contraction-side, that implicit sketch merging is the most stable source of efficiency, and that the approximation level acts as a reliable control knob for error.

The main contributions of this paper are:
\begin{itemize}
    \item We formulate approximate contraction for open-leg tensor networks under a full dense-materialization requirement, shifting the approximation target from the output tensor to the internal contraction process.
    \item We propose \method{}, which recursively maintains separator states on a partition tree, uses adaptive low-rank compression to control interface size, and employs implicit sketch-based merging to avoid explicitly forming inflated merged interfaces.
    \item We empirically validate the method on Ring-8, Tree-8, Grid $3\times 3$, and Random-8, showing a clear accuracy-time trade-off together with interpretable evidence about where the gains come from.
\end{itemize}
